% GENERATED FILE. Edit canonical lessons, then run npm run generate:books.
% canonical-source-hash: fnv1a64:63acae2d22179943
% canonical-lessons: TE-C44-look, TE-C44-listen, TE-C44-speak, TE-C44-write

\chapter{Verbs of Seeing and Saying}
\label{ch:te-see-say}
\hlchaptermodality{44}

\begin{chapteropening}
\textbf{By the end of this chapter:} \emph{I can tell someone to look, listen, speak and write in Telugu.}

Complete TE-C44-write and say all four verbs in order.
\end{chapteropening}

\section[cūḍu]{\te{చూడు} (cūḍu) — look}

\label{lesson:TE-C44-look}



\begin{quote}
Four verbs for taking things in and putting them out, one at a time.
\end{quote}

\subsection*{You'll want to know: \te{చూడు}}
\textbf{\te{చూడు}} (\emph{cūḍu}) — \textquotedblleft{}look\textquotedblright{}.

Also 'to see' and 'to look after', the same three-way stem as its Kannada cousin \emph{nōḍu} — the languages are close enough that the pattern travels.

\begin{grammarlens}[title={What you've built}]
One more. Three follow, and each reuses the ones before.
\end{grammarlens}

\subsection*{Your turn}
Say these aloud:

\begin{itemize}
\raggedright
  \item \emph{cūḍu}
  \item it once more, slowly
\end{itemize}

\begin{tcolorbox}[breakable,colback=teal!4,colframe=teal!35!black,title={Before you move on}]
What does \te{చూడు} mean? (\textquotedblleft{}look\textquotedblright{}.) Say it once more.
\end{tcolorbox}
\section[ālakiñcu]{\te{ఆలకించు} (ālakiñcu) — listen}

\label{lesson:TE-C44-listen}



\begin{quote}
Before the new one: what did \emph{cūḍu} mean?
\end{quote}

\subsection*{You'll want to know: \te{ఆలకించు}}
\textbf{\te{ఆలకించు}} (\emph{ālakiñcu}) — \textquotedblleft{}listen\textquotedblright{}.

'Listen attentively' — a step above simply hearing. Telugu keeps listening and asking separate, where Kannada merges them into one stem.

\begin{grammarlens}[title={What you've built}]
One more. Three follow, and each reuses the ones before.
\end{grammarlens}

\subsection*{Your turn}
Say these aloud:

\begin{itemize}
\raggedright
  \item \emph{ālakiñcu}
  \item it once more, slowly
  \item it after \emph{cūḍu}, so the two sit together
\end{itemize}

\begin{tcolorbox}[breakable,colback=teal!4,colframe=teal!35!black,title={Before you move on}]
What does \te{ఆలకించు} mean? (\textquotedblleft{}listen\textquotedblright{}.) Say it once more.
\end{tcolorbox}
\section[māṭlāḍu]{\te{మాట్లాడు} (māṭlāḍu) — speak}

\label{lesson:TE-C44-speak}



\begin{quote}
Before the new one: what did \emph{ālakiñcu} mean?
\end{quote}

\subsection*{You'll want to know: \te{మాట్లాడు}}
\textbf{\te{మాట్లాడు}} (\emph{māṭlāḍu}) — \textquotedblleft{}speak\textquotedblright{}.

From \emph{māṭa}, 'word', plus \emph{āḍu}, 'to play' — the same construction as Kannada's \emph{mātāḍu}, built from the same two pieces.

\begin{grammarlens}[title={What you've built}]
One more. Three follow, and each reuses the ones before.
\end{grammarlens}

\subsection*{Your turn}
Say these aloud:

\begin{itemize}
\raggedright
  \item \emph{māṭlāḍu}
  \item it once more, slowly
  \item it after \emph{ālakiñcu}, so the two sit together
\end{itemize}

\begin{tcolorbox}[breakable,colback=teal!4,colframe=teal!35!black,title={Before you move on}]
What does \te{మాట్లాడు} mean? (\textquotedblleft{}speak\textquotedblright{}.) Say it once more.
\end{tcolorbox}
\section[rāyi]{\te{రాయి} (rāyi) — write}

\label{lesson:TE-C44-write}



\begin{quote}
Before the new one: what did \emph{māṭlāḍu} mean?
\end{quote}

\subsection*{You'll want to know: \te{రాయి}}
\textbf{\te{రాయి}} (\emph{rāyi}) — \textquotedblleft{}write\textquotedblright{}.

The everyday 'write'. Telugu and Kannada each kept their own Dravidian stem here rather than borrowing a shared one.

\begin{grammarlens}[title={What you've built}]
That closes the run: look, listen, speak, write.
\end{grammarlens}

\subsection*{Your turn}
Say these aloud:

\begin{itemize}
\raggedright
  \item \emph{rāyi}
  \item it once more, slowly
  \item it after \emph{māṭlāḍu}, so the two sit together
\end{itemize}

\begin{tcolorbox}[breakable,colback=teal!4,colframe=teal!35!black,title={Before you move on}]
What does \te{రాయి} mean? (\textquotedblleft{}write\textquotedblright{}.) Say it once more.
\end{tcolorbox}
